\section{模型的建立与求解}

% 多问共享同一坐标系、状态方程或基础关系时，可先增加一个信息明确的共享模型二级标题。

\subsection{问题一的模型建立与求解}

\subsubsection{核心关系与模型建立}

% 定义对象、变量、目标与约束，给出本题关键关系、边界条件和必要推导。

\subsubsection{参数求取与数值求解}

% 写明输入、变换或离散、参数求取、迭代更新、约束处理、停止条件和报告精度。
% 国一模式：若本问包含三阶段以上、判断/循环/回退/早停、粗搜—细化、递归再生或状态转移，
% 必须在本问分析或求解段附近插入分问题流程图；图中变量、阈值、更新和停止条件须与正文及脚本一致。
% 每问至少插入一张解释图（机理/几何/流程）和一张结果或验证图，并在图前提出命题、图后定量解释。

\subsection{问题二的模型建立与求解}

% 按实际分问增删。单问较短时可不设三级标题，直接连续写模型与求解过程。
